% Documentação do compilador LastroLang (DIM0548), no modelo SBC de artigos.
% Gerar o PDF na raiz do projeto: make doc  ->  RELATORIO.pdf
% Cada seção fica em secoes/, um arquivo por issue do GitHub; o comentário no topo
% de cada arquivo diz o que entra nela.

\documentclass[12pt]{article}

\usepackage{sbc-template}        % modelo SBC: A4, margens, Times 12pt, títulos e legendas
\usepackage[T1]{fontenc}
\usepackage[utf8]{inputenc}
\usepackage[brazilian,shorthands=off]{babel}   % sem atalhos: " aparece em código
\usepackage{graphicx,url}
\usepackage{booktabs}            % tabelas sem linhas verticais, como pede o modelo
\usepackage{tabularx}
\usepackage{xcolor}
\usepackage{listings}
\usepackage[hidelinks,pdftitle={LastroLang: especificação da linguagem e implementação do compilador},
  pdfauthor={Caio Daniel Fonseca de Araújo, Flawbert Lorran da Silva Costa, Murilo Antonio Lima da Costa, Rodrigo Ferreira da Nobrega}]{hyperref}

% Estilo comum do relatório: código Lastro, código C e saída do compilador.
% Em tons de cinza, como recomenda o modelo SBC.

% Código curto no meio do texto e nas tabelas.
\newcommand{\cod}[1]{\texttt{#1}}

% Sem esticar o espaço entre parágrafos para encher a página: com tabelas e listagens,
% o \flushbottom do modelo abria buracos antes dos títulos.
\raggedbottom

\makeatletter
% O modelo só define \section e \subsection; este nível segue o mesmo padrão, em itálico.
\renewcommand\subsubsection{\@startsection{subsubsection}{3}{\z@}%
                       {-6\p@ \@plus -2\p@ \@minus -2\p@}%
                       {0.1\p@}%
                       {\normalfont\normalsize\bfseries\itshape}}

% O hyperref troca o \@lbibitem do modelo e volta a imprimir o rótulo [Autor Ano] na
% margem. Refaz a versão do modelo (referência sem rótulo), mantendo a âncora do link.
\def\@lbibitem[#1]#2{\item[]%
  \Hy@raisedlink{\hyper@anchorstart{cite.#2\@extra@b@citeb}\hyper@anchorend}%
  \if@filesw{\let\protect\noexpand\immediate
    \write\@auxout{\string\bibcite{#2}{#1}}}\fi\ignorespaces}
\makeatother

\lstdefinelanguage{Lastro}{
  morekeywords={const,proc,func,int,float,bool,if,else,while,for,to,
                break,continue,return,ref,true,false},
  sensitive=true,
  morecomment=[l]{//},
  morestring=[b]",
}

\lstset{
  basicstyle=\ttfamily\small,
  keywordstyle=\bfseries,
  commentstyle=\itshape\color{black!60},
  stringstyle=\color{black!75},
  columns=fullflexible,
  keepspaces=true,
  showstringspaces=false,
  upquote=true,
  frame=lines,
  framerule=0.4pt,
  rulecolor=\color{black!40},
  framesep=4pt,
  xleftmargin=0.8cm,
  xrightmargin=0.8cm,
  aboveskip=6pt,
  belowskip=6pt,
  % acentos em comentários e textos (listings lê byte a byte)
  literate={á}{{\'a}}1 {à}{{\`a}}1 {â}{{\^a}}1 {ã}{{\~a}}1
           {é}{{\'e}}1 {ê}{{\^e}}1 {í}{{\'i}}1 {ó}{{\'o}}1
           {ô}{{\^o}}1 {õ}{{\~o}}1 {ú}{{\'u}}1 {ç}{{\c{c}}}1
           {Á}{{\'A}}1 {É}{{\'E}}1 {Ó}{{\'O}}1 {Ç}{{\c{C}}}1,
}

\lstdefinestyle{lastro}{language=Lastro}
\lstdefinestyle{c}{language=C}
% saída do terminal: sem destaque de palavras
\lstdefinestyle{saida}{language={}, basicstyle=\ttfamily\small\color{black!85}}


\sloppy

\begin{document}

% Capa: logos no topo, título, autores e instituição centralizados e o resumo, com a
% disciplina e a data no rodapé. Usa as fontes do modelo SBC (\XVIPT no título).
\begin{titlepage}
  {\centering
  % alturas diferentes para as letras dos dois logos ficarem do mesmo tamanho (a UFRN
  % tem a legenda embaixo); centralizados na vertical e alinhados às margens do fio
  $\vcenter{\hbox{\includegraphics[height=1.8cm]{figuras/ufrn}}}$\hfill
  $\vcenter{\hbox{\includegraphics[height=1.4cm]{figuras/imd}}}$\par
  \vspace{8pt}
  {\color{black!35}\rule{\textwidth}{0.4pt}}\par
  \vspace{3.5cm}

  {\XVIPT\bfseries LastroLang: especificação da linguagem\\
    e implementação do compilador\par}
  \vspace{1.2cm}

  {\bfseries
    Caio Daniel Fonseca de Araújo\\
    Flawbert Lorran da Silva Costa\\
    Murilo Antonio Lima da Costa\\
    Rodrigo Ferreira da Nobrega\par}
  \vspace{0.8cm}

  Instituto Metrópole Digital (IMD)\\
  Universidade Federal do Rio Grande do Norte (UFRN)\\
  Natal, RN -- Brasil\par}
  \vspace{1.5cm}

  \begin{resumo}
    Este documento especifica a linguagem Lastro nos níveis léxico, sintático e
    semântico e descreve o compilador que a traduz para um C simplificado, aceito
    pelo \texttt{gcc} sem erros. Ele registra as principais decisões de
    implementação, as estruturas de dados e as verificações feitas sobre os
    programas, a organização do código-fonte, as instruções de geração e uso, as
    limitações da linguagem e da implementação e os pontos de melhoria futura.
  \end{resumo}

  \vfill
  {\centering
  {\color{black!35}\rule{\textwidth}{0.4pt}}\par
  \vspace{6pt}
  DIM0548 -- Engenharia de Linguagens -- Grupo 02\\
  1º de dezembro de 2026\par}
\end{titlepage}
\setcounter{page}{2}   % a capa conta como página 1 (sem número)

% Sumário em página própria. O modelo SBC suprime a numeração; o relatório numera as
% páginas a partir daqui para o sumário servir.
\pagestyle{plain}
\tableofcontents
\clearpage

% Issue #29 (Documentar a especificação da linguagem): parte de apresentação.
% - Visão geral da linguagem: para que serve, paradigma, tipagem estática
% - Um programa de exemplo completo em Lastro e o C gerado

\section{Apresentação da linguagem}
\label{sec:apresentacao}

% Item 1 da tarefa de submissão: especificação completa da linguagem.

\section{Especificação da linguagem}
\label{sec:especificacao}

% Issue #30 (Documentar a análise léxica), concluída; migrada do docs/analise-lexica.md.
% Issue #29: revisar contra o lexer.l final (#17 e #19 podem acrescentar tokens).

\subsection{Léxico}
\label{sec:lexico}

O analisador léxico de Lastro fica em \cod{src/lexer/lexer.l}, e o flex o
transforma na função \cod{yylex()} \cite{levine2009}. Quem comanda a leitura é o
analisador sintático, que pede um token por vez. A cada pedido, o léxico lê a
entrada padrão até fechar o próximo lexema e devolve o código do token. Se o
token tem valor, como um identificador ou um número, o lexema vai para
\cod{yylval.str}. A opção \cod{\%option yylineno} guarda a linha corrente, e
todas as mensagens de erro do compilador usam esse número, sejam elas léxicas,
sintáticas ou semânticas. O léxico só consulta a tabela de tipos, para devolver
\cod{TYPE\_ID} quando o identificador é o nome de um tipo. Ele entrega só o texto
do identificador, e o analisador sintático decide se aquilo é uma declaração ou
um uso, como fixamos no Problema~6.

\subsubsection*{Tokens e expressões regulares}

A especificação parte de três definições auxiliares: \cod{digito = [0-9]},
\cod{letra = [a-zA-Z]} e \cod{ident = \{letra\}(\{letra\}|\{digito\}|\_)*}.
A Tabela~\ref{tab:tokens} lista os tokens na mesma ordem em que as regras
aparecem no arquivo.

\begin{table}[ht]
\centering
\caption{Tokens de Lastro, na ordem das regras do \cod{lexer.l}}
\label{tab:tokens}
\small
\begin{tabularx}{\textwidth}{@{}l>{\raggedright\arraybackslash}X>{\raggedright\arraybackslash}X@{}}
\toprule
Classe & Lexema ou expressão regular & Token \\
\midrule
Palavra reservada & \cod{const proc func} & \cod{CONST PROC FUNC} \\
Palavra reservada & \cod{int float bool} & \cod{INT FLOAT BOOL} \\
Palavra reservada & \cod{if else while for to} & \cod{IF ELSE WHILE FOR TO} \\
Palavra reservada & \cod{switch case default} & \cod{SWITCH CASE DEFAULT} \\
Palavra reservada & \cod{break continue return ref} & \cod{BREAK CONTINUE RETURN REF} \\
Literal lógico & \cod{true false} & \cod{TRUE FALSE} \\
Atribuição & \cod{:=} & \cod{ASSIGN} \\
Relacional & \cod{= != <= < >= >} & \cod{EQ NEQ LE LT GE GT} \\
Aritmético & \cod{+ - * /} & \cod{PLUS MINUS STAR SLASH} \\
Lógico & \cod{\&\& || !} & \cod{AND OR NOT} \\
Tipo de retorno & \cod{->} & \cod{ARROW} \\
Delimitador & \cod{( ) [ ] \{ \}} & \cod{LPAREN RPAREN LBRACKET RBRACKET LBRACE RBRACE} \\
Separador & \cod{, ; .} & \cod{COMMA SEMI DOT} \\
Literal de texto & \cod{\textbackslash"[\textasciicircum\textbackslash"\textbackslash\textbackslash\textbackslash n]*\textbackslash"} & \cod{LITERAL\_TEXTO} \\
Texto sem fechar & \cod{\textbackslash"[\textasciicircum\textbackslash"\textbackslash\textbackslash\textbackslash n]*} & \cod{ERRO} \\
Literal real & \cod{\{digito\}+"."\{digito\}+} & \cod{LITERAL\_FLOAT} \\
Literal inteiro & \cod{\{digito\}+} & \cod{LITERAL\_INT} \\
Identificador & \cod{\{ident\}} & \cod{ID} ou \cod{TYPE\_ID} \\
Erro & \cod{.} (qualquer outro caractere) & \cod{ERRO} \\
\bottomrule
\end{tabularx}
\end{table}

A linguagem tem 20 palavras reservadas, todas em minúsculas. Como ela diferencia
maiúsculas de minúsculas, \cod{If} e \cod{WHILE} são identificadores comuns.
\cod{ler} e \cod{escrever} também são identificadores: chegam ao analisador
sintático como \cod{ID} e entram na tabela de símbolos como subprogramas
embutidos, de modo que a análise semântica confere uma chamada a eles do mesmo
jeito que confere qualquer outra.

Todo identificador começa com uma letra ASCII e continua com letras, dígitos ou
sublinhado. Deixamos as letras acentuadas de fora porque o código de sistemas
financeiros conversa com bibliotecas C e com sistemas legados (Problema~1), e um
nome como \cod{preço} não teria equivalente direto no C gerado. O literal
inteiro é uma sequência de dígitos, sem sinal. Em \cod{-5}, o \cod{-} é o
operador unário da gramática. O literal real exige dígitos dos dois lados do
ponto (\cod{2.5}; \cod{2.} e \cod{.5} não são aceitos). Em \cod{2.5}, as regras
de inteiro e de real casam ao mesmo tempo, e o casamento mais longo escolhe
\cod{LITERAL\_FLOAT}. Valores monetários continuam sendo guardados como inteiros
(em centavos ou pontos-base), como decidido no Problema~1. O \cod{float} existe
para os cálculos não monetários e para a coerção de \cod{int} para \cod{float}
exigida na avaliação.

O literal de texto fica entre aspas, na mesma linha e sem \cod{\textbackslash},
e vai para o C exatamente como foi escrito. A linguagem não tem tipo
\cod{string}: o texto só é aceito como argumento de \cod{escrever}, para as
mensagens exigidas nos problemas de avaliação. Quando as aspas não fecham na
linha, ou quando aparece uma \cod{\textbackslash}, só a regra de texto sem
fechar casa, e o léxico reporta \cod{texto sem aspas de fechamento}.

O léxico já reconhece \cod{SWITCH}, \cod{CASE}, \cod{DEFAULT} e \cod{DOT}, mas a
gramática ainda não usa esses tokens em nenhuma produção. Eles ficam reservados
para o comando \cod{switch} e para o acesso a campo de registro, e o yacc os
aponta como terminais não usados. Quando um deles aparece num programa, o léxico
o aceita e o analisador sintático acusa erro de sintaxe.

\subsubsection*{Casamento mais longo e palavras reservadas antes de identificador}

O flex desempata com duas regras \cite{aho2007}. Pela primeira, o casamento mais
longo (\emph{longest match}), vence a expressão que consome mais caracteres a
partir da posição atual. Se duas expressões casam o mesmo número de caracteres,
vale a segunda regra: ganha a que foi declarada primeiro no arquivo.

Com o casamento mais longo, os operadores que começam igual não precisam de
tratamento especial. Em \cod{<=}, tanto \cod{"<="} quanto \cod{"<"} casam, mas a
primeira leva dois caracteres e ganha. O mesmo vale para \cod{:=} e \cod{=},
\cod{>=} e \cod{>}, \cod{!=} e \cod{!}, \cod{->} e \cod{-}, e por isso a ordem
dessas regras no arquivo não muda nada. O comentário segue a mesma lógica:
\cod{"//".*} sempre consome mais que \cod{"/"}, então duas barras seguidas nunca
viram duas divisões.

Nas palavras reservadas, quem decide é a ordem. \cod{while} casa \cod{"while"} e
também \cod{\{ident\}}, com o mesmo comprimento, e só é lido como palavra
reservada porque essas regras estão antes da regra de identificador no arquivo.
Se a ordem fosse invertida, o léxico nunca devolveria uma palavra reservada. O
casamento mais longo cobre o caso inverso. Em \cod{ifx}, a regra \cod{"if"} casa
dois caracteres e \cod{\{ident\}} casa três, então \cod{ifx} é um identificador,
e o programador pode dar a uma variável um nome que começa com uma palavra
reservada.

\subsubsection*{Comentários, espaços e erro léxico}

O léxico reconhece comentários e espaços e os descarta sem gerar token. O
comentário vai de \cod{//} até o fim da linha (\cod{"//".*}), porque o \cod{.}
do flex não casa a quebra de linha. A linguagem não tem comentário de bloco,
então ninguém esquece um comentário aberto pelo caminho, e cada linha pode ser
lida sozinha, o que atende à prioridade de legibilidade do Problema~1. A
expressão \cod{[ \textbackslash t\textbackslash r\textbackslash n]+} cobre os
espaços, que só separam lexemas. O \cod{\textbackslash r} está ali para aceitar
arquivos com quebra de linha do Windows. As quebras de linha não geram token,
mas o \cod{yylineno} as conta.

A última regra, \cod{.}, pega qualquer caractere que nenhuma outra regra
aceitou: \cod{@}, \cod{\#}, um \cod{:} sozinho, um \cod{\&} ou \cod{|} sozinho,
ou cada byte de uma letra acentuada. A ação chama
\cod{yyerror("caractere inválido")}, que imprime a linha e o lexema, e devolve o
token \cod{ERRO}. Nenhuma produção usa esse token, então o analisador sintático
rejeita o programa no mesmo ponto e o compilador sai com código 1. Com
\cod{tests/invalido/05\_caractere\_invalido.txt}, a saída é:

\begin{lstlisting}[style=saida]
linha 3: caractere inválido perto de '@'
linha 3: syntax error perto de '@'
REJEITADO: programa contém erro
\end{lstlisting}

Dois dos erros previstos no Problema~6 formam sequências de tokens válidas e por
isso passam pelo léxico. Quem os detecta é o analisador sintático. O literal
fracionário incompleto \cod{10.} vira \cod{LITERAL\_INT DOT}, e a mensagem é
\cod{syntax error perto de '.'}. O literal mal formado \cod{3x} vira
\cod{LITERAL\_INT ID}, com a mensagem \cod{syntax error perto de 'x'}. Nos dois
casos o programa é rejeitado na linha certa.

% Issue #29: nível sintático.
% - Gramática BNF completa, gerada do parser.y final (também em docs/gramatica-bnf.md)
% - Precedência e associatividade dos operadores
% - Chaves obrigatórias eliminando o dangling-else

\subsection{Sintaxe}
\label{sec:sintaxe}

% Issue #29: nível semântico.
% - Tipos (int, float, bool, arrays, matrizes, registros) e equivalência de tipos
% - Coerção de int para float e o que é proibido
% - Escopos: declarar antes de usar, duplicata no mesmo escopo, shadowing em escopo aninhado
% - Constantes, subprogramas (proc/func, valor e ref, recursão) e os embutidos ler/escrever

\subsection{Semântica}
\label{sec:semantica}


% Item 2 da tarefa de submissão: decisões de implementação, estruturas de dados,
% verificações e organização dos códigos-fonte.

\section{Implementação do compilador}
\label{sec:implementacao}

% Issue #52 (Documentar a organização do código e as decisões de implementação).
% - Diretórios e conteúdo de cada arquivo (src/lexer, src/parser, src/semantic, src/gen,
%   tests, programas, Makefile)
% - Figura do fluxo: entrada -> yylex() -> yyparse() com ações semânticas -> C -> gcc
% - Passada única sem árvore sintática e o atributo Atr (código C, tipo, linha, eh_const)
% - C simplificado: if (!c) goto, rótulos _Ln, temporários _tN; bool -> int, float -> double,
%   const -> static const, ref -> ponteiro, array por valor com memcpy
% - Embutidos ler/escrever com uma função C por tipo
% - Saída e erros: stdout/stderr, código de saída, opções -t e -s
% - Decisões tomadas e o motivo

\subsection{Organização do código e decisões de implementação}
\label{sec:organizacao}

% Issue #31 (Documentar a tabela de símbolos e os escopos).
% - As três tabelas (símbolos escopo#nome, subprogramas, tipos), a tabela hash e o TypeDesc
% - table_insert/lookup, sub_insert/lookup, type_insert/lookup, scope_push/pop/level,
%   find_name e nome_livre
% - Shadowing: nome_livre olha só o topo; find_name acha o mais interno
% - Exemplo do estado das tabelas e da pilha (saída da opção -s)

\subsection{Tabela de símbolos e escopos}
\label{sec:tabelas}

% Issue #32 (Documentar condicionais e repetição).
% - if/else if/else, while e for ... to (inclusivo)
% - Dangling-else eliminado pelas chaves obrigatórias
% - break/continue e os rótulos do laço mais interno
% - Escopo do for e tradução para C com if (!c) goto, com exemplo

\subsection{Condicionais e repetição}
\label{sec:controle}

% Issue #33 (Documentar os subprogramas).
% - proc x func; cabeçalho registrado antes do corpo (recursão)
% - Parâmetros por valor e ref (ponteiro; array por valor com memcpy; coerção int -> float)
% - Retorno, verificação das chamadas e embutidos ler/escrever (uma função C por tipo)

\subsection{Subprogramas}
\label{sec:subprogramas}

% Issue #34 (Documentar as verificações semânticas).
% - Tabela única: regra, mensagem de erro e exemplo de cada verificação
% - Tipos, coerções, faixas, duplicidade e shadowing, declaração, constantes, chamadas,
%   texto só em escrever e controle de fluxo

\subsection{Verificações realizadas}
\label{sec:verificacoes}


% Issue #35 (Documentar as instruções de geração e uso e submeter o PDF final).
% - Pré-requisitos (gcc, flex, yacc/bison, TeX Live para o relatório)
% - make, compilar um programa, gerar o C, rodar, make test, make doc
% - Opções -t e -s, mensagens de erro e códigos de saída

\section{Instruções de geração e uso}
\label{sec:uso}

% Issue #53 (Documentar as limitações e as melhorias futuras): limitações.
% Falta acrescentar:
% - Linguagem: sem tipo string (texto só em escrever), sem ponteiros nem alocação dinâmica,
%   sem switch, literal real só no formato 2.5, coerção só de int para float,
%   tamanho de array fixo em compilação, e o que não sair até a entrega
% - Implementação: passada única sem árvore sintática; limites fixos de 100 níveis de
%   escopo (MAX_SCOPE_DEPTH) e 256 laços aninhados (MAX_LACOS); faixa de índice só no que a #12 cobrir

\section{Limitações}
\label{sec:limitacoes}

O analisador sintático usa a recuperação de erro padrão do yacc e para no
primeiro erro de sintaxe. Como todo \cod{ERRO} provoca um erro de sintaxe, o
compilador relata só o primeiro caractere inválido do programa. O Problema~6
previa que o léxico seguisse adiante depois de um erro, e a implementação atual
não faz isso. O léxico também não confere a faixa do literal inteiro:
\cod{99999999999999999999} passa e só daria problema ao compilar o C gerado.

% Issue #53 (Documentar as limitações e as melhorias futuras): melhorias.
% - Árvore sintática e uma passada de otimização
% - Recuperação de erros sintáticos para reportar mais de um erro por compilação
% - switch, potência, tipo string, ponteiros e alocação dinâmica

\section{Melhorias futuras}
\label{sec:melhorias}


% Lista as referências mesmo sem \cite (as seções ainda estão sendo escritas).
\nocite{*}
\bibliographystyle{sbc}
\bibliography{referencias}

\end{document}
